\section{The Methods}

\begin{frame}{One Engine, Five Methods}
  \textbf{Shared core:} every method runs on the same incremental machinery ---
  per-move crossing deltas, spatial grids for candidate lookup, strict validity
  (no vertex on a foreign edge), and best-snapshot bookkeeping.
  \medskip
  \begin{center}
    \footnotesize
    \begin{tabular}{llll}
      \toprule
      \textbf{Method} & \textbf{Pipeline} & \textbf{Escapes local optima by} & \textbf{Binary} \\
      \midrule
      \texttt{sa}        & two-phase SA                    & temperature waves            & \texttt{sakgd} \\
      \texttt{sa-stress} & \textbf{stress init} $\to$ SA   & better start + waves         & \texttt{sfdp} + \texttt{sakgd} \\
      \texttt{staged}    & LNS (30\%) $\to$ SA (70\%)      & SA stage after greedy LNS    & \texttt{approach1} \\
      \texttt{staged-adaptive} & adaptive-LNS $\to$ SA     & growing neighbourhoods       & \texttt{approach1} \\
      \texttt{ils}       & SA rounds + random kicks        & perturbation between rounds  & \texttt{approach1} \\
      \bottomrule
    \end{tabular}
  \end{center}
  \medskip
  \textbf{Two-phase SA} (after the winning GD'25 entry):
  phase 1 minimises total crossings $\chi$ (untangles the drawing),
  phase 2 minimises the bottleneck $k$ with a dual fitness:
  $\Delta E = \Delta K_{\text{local}}$ if a touched edge's $k$ changes,
  else $\Delta\chi/\chi$ as a tiny tie-break.
\end{frame}

\begin{frame}{Inside the SA Loop: Moves, Waves, Snapshots}
  \begin{columns}[T]
    \begin{column}{0.52\textwidth}
      \textbf{One move:}
      \begin{enumerate}
        \item pick a vertex --- crossing-weighted; in phase 2, biased to
              \emph{$k$-critical} endpoints,
        \item propose a position (Gaussian around current; radius shrinks
              with temperature),
        \item reject if occupied or vertex-on-edge,
        \item Metropolis: accept if $\Delta E\le 0$, else with
              probability $e^{-\Delta E/T}$.
      \end{enumerate}
      \medskip
      \textbf{Temperature waves:} inner loop cools $T$ per move;
      the outer loop restarts each wave from the best-so-far solution
      at a slightly lower starting temperature.
    \end{column}
    \begin{column}{0.45\textwidth}
      \textbf{Best-snapshot bookkeeping:}
      \begin{itemize}
        \item \texttt{saveBest()} --- whenever $(k,\chi)$ improves, copy the
              positions into \texttt{bestPos} (three fields, \emph{in RAM}).
        \item \texttt{restoreBest()} --- copy \texttt{bestPos} back and rebuild
              the tables: at wave ends, and once more before writing output.
      \end{itemize}
      \medskip
      \textbf{Stateless by design:} no files, no caches, nothing survives the
      process. Two runs on the same input are fully independent ---
      warm-starting across runs exists only as an explicit, off-by-default
      pipeline flag.
    \end{column}
  \end{columns}
\end{frame}

\begin{frame}{Method Shoot-Out: Same Seeds, Same Budgets}
  All five methods, all nine instances, one controlled run
  (4 workers $\times$ seed 42--45, budgets 2/4/6\,min by size; best of 4 workers shown):
  \begin{center}
    \footnotesize
    % RESULTS-TABLE-PLACEHOLDER (filled from results/runs/<final-method-comparison>)
    \begin{tabular}{lrrrrr}
      \toprule
      \textbf{Graph} & \texttt{sa} & \texttt{sa-stress} & \texttt{ils} & \texttt{staged} & \texttt{staged-adaptive} \\
      \midrule
      Automatic-1 & & & & & \\
      Automatic-9 & & & & & \\
      \bottomrule
    \end{tabular}
  \end{center}
\end{frame}
